\textit{(                                               HIBA)}

\begin{lstlisting}
# 1. Introduction

## 1.1 Motivation

Standard scaled dot-product attention (Vaswani et al., 2017) computes the attention matrix at $\mathcal{O}(L^2)$ time and memory complexity. While sparse attention patterns (such as Longformer and BigBird) reduce this to $\mathcal{O}(L \sqrt{L})$ or $\mathcal{O}(L)$, they permanently sever global token-to-token routing pathways. This strict sparsity often causes catastrophic accuracy degradation when injected post-training into dense models. 

## 1.2 Problem Statement

Given a pretrained transformer $\mathcal{M}$, we seek a modified attention function $\mathcal{A}'$ that replaces $\mathcal{A}$ such that memory scales sub-quadratically, yet global context is not permanently lost. We impose the constraints:
1. **No retraining of base weights.** 
2. **Robust architectural support.** Must seamlessly handle nested encoders (RoBERTa), fused weights (GPT-2), and modern features like GQA and RoPE without double-projection errors.

## 1.3 Our Approach

We propose **Hierarchical Intra-Block and Inter-Block Attention (HIBA)**. HIBA processes tokens in local blocks (e.g., $B=64$) for precise local semantics. For global context, it sparsely attends to a dynamic Top-K subset of tokens from neighboring blocks. Most critically, to prevent the "lost in the middle" degradation inherent to sparse architectures, HIBA employs **Periodic Dense Rectification**, executing a full $\mathcal{O}(L^2)$ attention pass at a configurable interval (e.g., every 4th layer). 

\end{lstlisting}

